\appsection{Hyperprior Derivation by Distributional Family}
\label{app:hyperprior_families}
%==============================================================================

This appendix provides the complete derivation of hyperprior parameters from
$(\rho, \kappa)$ for three distributional families covering the most common
support types in applied work: the Normal-Inverse-Gamma (support on
$\mathbb{R}$), Gamma-Inverse-Gamma (support on $(0,\infty)$), and Beta-Beta
(support on $[0,1]$). Each subsection presents the hierarchy, the closed-form
translation from researcher inputs to hyperprior parameters, the feasibility
region in $(\rho, \kappa)$ space, and graphical calibrations.

%----------------------------------------------------------------------
\subsection{Normal-Inverse-Gamma (support on $\mathbb{R}$)}
\label{app:nig}

Consider the standard Normal-Inverse-Gamma hierarchical model:
\begin{align}
\bv_i \mid \mu, \sigma^2 &\sim N(\mu, \sigma^2) \\
\mu \mid \mu_0, \tau_0^2 &\sim N(\mu_0, \tau_0^2) \\
\sigma^2 \mid \alpha_0, \beta_0 &\sim \text{IG}(\alpha_0, \beta_0)
\end{align}
Here $\hypem = \mu$ and $\hypev = \sigma^2$, so the hyperprior parameters are
$\hprior = (\mu_0, \tau_0^2, \alpha_0, \beta_0)$.

From (\ref{eq:hyperm_solved}), the mean hyperprior parameters follow
immediately, since the Normal mean and variance are the distribution's
parameters:
\begin{align}
  \mu_0 &= \dpm \\
  \tau_0^2 &= \rho\,\dpv
\end{align}
The conditional variance is the parameter $\sigma^2$ itself, so $\kappa =
\text{SD}[\sigma^2] / \E{\sigma^2}$ maps directly to the Inverse-Gamma
hyperprior. Using $\E{\sigma^2} = \beta_0 / (\alpha_0 - 1)$
and $\Var{\sigma^2} = \beta_0^2 / [(\alpha_0 - 1)^2(\alpha_0 - 2)]$,
equations (\ref{eq:hyperv_solved}) yield:
\begin{align}
  \alpha_0 &= 2 + \frac{1}{\kappa^2} \\
  \beta_0 &= (1 - \rho)\,\dpv\,(\alpha_0 - 1)
\end{align}
The shape parameter $\alpha_0$ is determined solely by $\kappa$, with
$\beta_0$ then following from $\rho$ and $\dpv$.

The NIG family admits the full region $\rho \in (0,1)$, $\kappa \in
(0,\infty)$. The Inverse-Gamma requires $\alpha_0 > 2$ for the variance of
$\sigma^2$ to exist, which holds for all $\kappa > 0$ since $\alpha_0 = 2 +
1/\kappa^2 > 2$ automatically. There are no joint constraints linking $\rho$
and $\kappa$.

\begin{figure}[htbp]
  \centering
  \small (a) Between-population SD: $\text{SD}[\mu] = \tau_0$. \\
  \includegraphics[width=0.6\textwidth]{../plots/paper/nig_sensitivity_between.pdf}\\[-.75em]

  \vspace{0.8em}
  \small (b) Within-population SD: $\E{\sigma}$. \\
  \includegraphics[width=0.6\textwidth]{../plots/paper/nig_sensitivity_within.pdf}\\[-.75em]

  \vspace{0.8em}
  \small (c) Within-population SD IQR. \\
  \includegraphics[width=0.6\textwidth]{../plots/paper/nig_sensitivity_iqr.pdf}\\[-.75em]

  \caption{Hierarchical variance allocation for the NIG model across
    $(\rho, \kappa)$, with $\dpm = 0$ and $\dpv = 1$.}
  \label{fig:nig_sensitivity}
\end{figure}

Figure~\ref{fig:nig_sensitivity} displays the hierarchical variance allocation
across the $(\rho, \kappa)$ grid, setting $\dpm = 0$ and $\dpv = 1$. Each
$(\rho, \kappa)$ pair uniquely determines the hyperprior parameters $(\mu_0,
\tau_0^2, \alpha_0, \beta_0)$ through the moment conditions
(\ref{eq:hyperm_solved})--(\ref{eq:hyperv_solved}). Panel~(a) shows that the
between-population SD depends only on $\rho$, as the horizontal color bands
confirm: $\kappa$ plays no role. This follows directly from the moment
condition $\tau_0 = \sqrt{\rho\,\dpv}$. Panel~(b) shows the within-population
SD, $\E{\sigma}$: higher $\rho$ allocates more variance to population-mean
uncertainty, leaving less for heterogeneity, so the within-population SD
decreases as $\rho$ increases. Panel~(c) reveals the role of $\kappa$. At low
$\kappa$ the within-population SD is tightly concentrated (the IQR is small),
reflecting near-certainty about the level of heterogeneity. As $\kappa$ grows
the IQR increases, reflecting greater uncertainty about how heterogeneous the
population is.

\begin{figure}[htbp]
  \centering
  \small (a) KS distance: $\sup_v |F_{\text{compound}}(v) - F_{\text{direct}}(v)|$. \\
  \includegraphics[width=0.6\textwidth]{../plots/paper/nig_ks_distance.pdf}\\[-.75em]

  \vspace{0.8em}
  \small (b) CDF difference at four $(\rho, \kappa)$ points. \\
  \includegraphics[width=0.8\textwidth]{../plots/paper/nig_cdf_diff_curves.pdf}\\[-.75em]

  \caption{Distributional approximation quality for the NIG hierarchy,
    $\dpm = 0$, $\dpv = 1$.}
  \label{fig:nig_distributional}
\end{figure}

Figure~\ref{fig:nig_distributional} quantifies the distributional discrepancy
between the hierarchical predictive and the direct prior. Panel~(a) shows the
Kolmogorov--Smirnov distance across the full $(\rho, \kappa)$ grid. For moderate
values such as those used in the application ($\rho = 0.2$, $\kappa =
0.25$), the supremum CDF difference is approximately 0.3 percentage points.
Even at the extreme corner ($\rho \to 0$, $\kappa = 1.7$), it remains below
5~percentage points. Panel~(b) reveals where the discrepancy arises: the
compound distribution is flatter-peaked and heavier-tailed than the direct
prior, consistent with the scale-mixture structure of the hierarchy. The CDF
difference is negative in the shoulders ($|v| \lesssim 1.8\sigma^*$), where the
compound places less mass, and positive in the tails ($|v| \gtrsim
1.8\sigma^*$), where it places more. All four curves share this qualitative
shape; only the amplitude varies with $\rho$ and $\kappa$.

The first two moments of the marginal prior on individual quantities are
invariant across the grid by construction; only the hierarchical allocation
changes. The specific numerical values are family-dependent, but these
qualitative patterns generally hold for any distributional choice satisfying the
moment conditions.

%----------------------------------------------------------------------
\subsection{Gamma-Inverse-Gamma (support on $(0, \infty)$)}
\label{app:gamma_ig}

For positive-valued quantities such as income, prices, or time
preferences, the Gamma-Inverse-Gamma hierarchy provides a natural model:
\begin{align}
  v_i \mid \mu, \sigma^2
    &\sim \text{Gamma}\!\left(
      \frac{\mu^2}{\sigma^2},\;
      \frac{\mu}{\sigma^2}
    \right) \\
  \mu \mid a_\mu, b_\mu &\sim \text{Gamma}(a_\mu, b_\mu) \\
  \sigma^2 \mid \alpha_0, \beta_0 &\sim \text{IG}(\alpha_0, \beta_0)
\end{align}
where the individual-level Gamma is parameterized by shape and rate so that
$\E{v_i \mid \mu, \sigma^2} = \mu$ and $\Var{v_i \mid \mu, \sigma^2} =
\sigma^2$. The conditional variance is the parameter $\sigma^2$ itself, as in
the NIG case.

The mean hyperprior uses $\E{\mu} = \dpm$ and $\Var{\mu} = \rho\,\dpv$,
giving shape $a_\mu$ and rate $b_\mu$:
\begin{align}
  a_\mu &= \frac{(\dpm)^2}{\rho\,\dpv} \\
  b_\mu &= \frac{\dpm}{\rho\,\dpv}
\end{align}
Because $\sigma^2$ is a parameter with its own independent hyperprior, the
$\kappa$ translation is identical to the NIG case:
\begin{align}
  \alpha_0 &= 2 + \frac{1}{\kappa^2} \\
  \beta_0 &= (1 - \rho)\,\dpv\,(\alpha_0 - 1)
\end{align}

As with NIG, the Inverse-Gamma variance hyperprior admits all $\kappa > 0$.
The Gamma mean hyperprior requires $a_\mu > 0$ and $b_\mu > 0$, which hold
whenever $\dpm > 0$ and $\rho\,\dpv > 0$. The feasibility region is therefore
the full rectangle $\rho \in (0,1)$, $\kappa \in (0, \infty)$, the same as
NIG.

The within-population statistics (within-population SD and its IQR) depend only
on the Inverse-Gamma hyperprior with $\alpha_0 = 2 + 1/\kappa^2$ and $\beta_0
= (1-\rho)\,\dpv\,(\alpha_0 - 1)$. These are identical to the NIG case, so
the sensitivity heatmaps in Figure~\ref{fig:nig_sensitivity} apply without
modification: panel~(a) gives the between-population SD, panel~(b) the
within-population SD, and panel~(c) the within-population SD IQR for the
Gamma-IG hierarchy as well. The between-population SD is
$\sqrt{\rho\,\dpv}$ regardless of family.

The two families differ in the shape of the compound distribution:
the NIG compound is symmetric with Student-$t$-like tails, while the
Gamma-Inverse-Gamma compound is right-skewed.

\begin{figure}[htbp]
  \centering
  \small (a) KS distance: $\sup_v |F_{\text{compound}}(v) - F_{\text{direct}}(v)|$. \\
  \includegraphics[width=0.6\textwidth]{../plots/paper/gamma_ks_distance.pdf}\\[-.75em]

  \vspace{0.8em}
  \small (b) CDF difference at four $(\rho, \kappa)$ points. \\
  \includegraphics[width=0.8\textwidth]{../plots/paper/gamma_cdf_diff_curves.pdf}\\[-.75em]

  \caption{Distributional approximation quality for the Gamma-IG hierarchy,
    $\dpm = 1$, $\dpv = 1$.}
  \label{fig:gamma_distributional}
\end{figure}

Figure~\ref{fig:gamma_distributional} shows the distributional approximation
quality for the Gamma-Inverse-Gamma case. Unlike the NIG, the Gamma-IG
compound has a nonzero KS distance even at small $(\rho, \kappa)$: the minimum
across the grid is approximately 0.005, reflecting an irreducible
distributional mismatch between the Gamma compound and the direct Gamma prior
that persists even as the hyperprior concentrates. The CDF difference curves
are no longer antisymmetric, as the right-skewed compound distribution differs
from the direct Gamma prior in a pattern that depends on the skewness of both
distributions.

%----------------------------------------------------------------------
\subsection{Beta-Beta (support on $[0,1]$)}
\label{app:beta_beta}

For bounded quantities normalized to $(0,1)$, the Beta-Beta hierarchy uses a
mean-dispersion parameterization. Let $z_i^k$ denote the normalized value of
quantity $k$ for individual $i$:
\begin{align}
  z_i^k \mid \mu_k, \psi_k
    &\sim \text{Beta}\!\left(
      \mu_k \tfrac{1 - \psi_k}{\psi_k},\;
      (1 - \mu_k) \tfrac{1 - \psi_k}{\psi_k}
    \right) \label{eq:app_bb_individual} \\
  \mu_k &\sim \text{Beta}(a_\mu^k, b_\mu^k)
    \label{eq:app_bb_mu} \\
  \psi_k &\sim \text{Beta}(a_\psi^k, b_\psi^k)
    \label{eq:app_bb_psi}
\end{align}
where $\mu_k \in (0,1)$ governs the population mean and $\psi_k \in (0,1)$
governs dispersion. The conditional moments are $\E{z_i^k \mid \mu_k, \psi_k}
= \mu_k$ and $\Var{z_i^k \mid \mu_k, \psi_k} = \mu_k(1 - \mu_k)\psi_k$.

The mean-dispersion parameterization is convenient because it allows
independent hyperpriors on $\mu_k$ and $\psi_k$. An alternative would be to
parameterize the individual-level Beta in terms of its mean and variance
directly, but $\Var{z \mid \mu, \sigma^2} = \sigma^2$ would require $\sigma^2 <
\mu(1-\mu)$, creating a conditional domain that precludes independent
hyperpriors. The normalization $\psi_k = \sigma_k^2 / [\mu_k(1 - \mu_k)]$
removes this dependence.

The law of total variance decomposes the direct prior variance:
\begin{align}
  \label{eq:app_bb_total_var}
  \dpv_k &= \Var{\mu_k} + \E{\mu_k(1 - \mu_k)\psi_k} \\
  \Var{\mu_k} &= \rho_k\,\dpv_k \\
  \E{\mu_k(1 - \mu_k)\psi_k} &= (1 - \rho_k)\,\dpv_k
\end{align}
Matching $\E{\mu_k} = \dpm_k$ and $\Var{\mu_k} = \rho_k\,\dpv_k$
by standard Beta moment-matching, the mean hyperprior parameters are:
\begin{align}
  a_\mu^k &= \dpm_k
    \left(\frac{\dpm_k(1 - \dpm_k)}{\rho_k\,\dpv_k} - 1\right) \\
  b_\mu^k &= (1 - \dpm_k)
    \left(\frac{\dpm_k(1 - \dpm_k)}{\rho_k\,\dpv_k} - 1\right)
\end{align}
Using the independence of $\mu_k$ and $\psi_k$, the within-population
term determines $\E{\psi_k}$:
\begin{equation}
  \E{\psi_k} = \frac{(1 - \rho_k)\,\dpv_k}{\E{\mu_k(1 - \mu_k)}}
\end{equation}
where $\E{\mu_k(1 - \mu_k)} = \dpm_k(1-\dpm_k) - \rho_k\,\dpv_k$ is computed
from the already-determined Beta$(a_\mu^k, b_\mu^k)$ hyperprior.

The parameter $\kappa_k$ is the coefficient of variation of the conditional
variance $\mu_k(1-\mu_k)\psi_k$, as defined in Section
\ref{sec:variance_decomposition}. Because the conditional variance is a product
of two independent random variables, its variance involves cross-moments:
\begin{equation}
  \Var{\mu_k(1 - \mu_k)\psi_k}
  = \E{\mu_k^2(1 - \mu_k)^2}\,\Var{\psi_k}
  + \Var{\mu_k(1 - \mu_k)}\,\E{\psi_k}^2
\end{equation}
Setting the left-hand side to $\kappa_k^2 (1-\rho_k)^2 (\dpv_k)^2$ and solving
for $\Var{\psi_k}$:
\begin{equation}
  \label{eq:app_bb_var_psi}
  \Var{\psi_k}
  = \frac{\kappa_k^2\,(1-\rho_k)^2\,(\dpv_k)^2
          - \Var{\mu_k(1-\mu_k)}\,\E{\psi_k}^2}
         {\E{\mu_k^2(1-\mu_k)^2}}
\end{equation}
The required moments are closed-form functions of $(a_\mu^k, b_\mu^k)$.
For $\mu_k \sim \text{Beta}(a,b)$ with $S = a + b$:
\begin{align}
  \E{\mu_k(1-\mu_k)} &= \frac{ab}{S(S+1)} \\
  \E{\mu_k^2(1-\mu_k)^2} &= \frac{a(a+1)b(b+1)}{S(S+1)(S+2)(S+3)} \\
  \Var{\mu_k(1-\mu_k)} &= \E{\mu_k^2(1-\mu_k)^2} - \E{\mu_k(1-\mu_k)}^2
\end{align}
Standard Beta moment-matching then gives the dispersion hyperprior parameters:
\begin{align}
  \nu_\psi^k &= \frac{\E{\psi_k}(1 - \E{\psi_k})}{\Var{\psi_k}} - 1 \\
  a_\psi^k &= \E{\psi_k}\,\nu_\psi^k \\
  b_\psi^k &= (1 - \E{\psi_k})\,\nu_\psi^k
\end{align}

The Beta-Beta family imposes two constraints on $(\rho, \kappa)$.
First, the mean hyperprior requires $\rho\,\dpv_k < \dpm_k(1-\dpm_k)$, the
Bernoulli bound on the variance of a $[0,1]$-valued random variable. This
constrains $\rho$ given the direct prior.
Second, $\Var{\psi_k} > 0$ in (\ref{eq:app_bb_var_psi}) requires:
\begin{equation}
  \kappa > \frac{\sqrt{\Var{\mu_k(1-\mu_k)}}}{\E{\mu_k(1-\mu_k)}}
  = \text{CV}[\mu_k(1-\mu_k)]
\end{equation}
This lower bound on $\kappa$ reflects the irreducible variation in conditional
variance arising from $\mu_k$ alone: even with $\psi_k$ fixed, the conditional
variance $\mu_k(1-\mu_k)\psi_k$ varies because $\mu_k$ is random. The bound
increases with $\rho$ (which increases $\Var{\mu_k}$) and is largest when
$\dpm_k$ is near the boundary of $[0,1]$.
Additionally, $\Var{\psi_k} < \E{\psi_k}(1-\E{\psi_k})$ is required for the
Beta dispersion hyperprior to be proper, imposing an upper bound on $\kappa$
for each $(\dpm_k, \dpv_k, \rho_k)$ configuration.
The feasibility region is therefore a proper subset of $(0,1) \times
(0,\infty)$, with both lower and upper bounds on $\kappa$ that depend on
$\rho$, $\dpm_k$, and $\dpv_k$. For typical prior configurations with
$\dpm_k$ not too close to 0 or 1, the feasible range of $\kappa$ is wide.

\begin{figure}[htbp]
  \centering
  \small (a) Between-population SD. \\
  \includegraphics[width=0.6\textwidth]{../plots/paper/bb_sensitivity_between.pdf}\\[-.75em]

  \vspace{0.8em}
  \small (b) Within-population SD. \\
  \includegraphics[width=0.6\textwidth]{../plots/paper/bb_sensitivity_within.pdf}\\[-.75em]

  \vspace{0.8em}
  \small (c) Within-population SD IQR. \\
  \includegraphics[width=0.6\textwidth]{../plots/paper/bb_sensitivity_iqr.pdf}\\[-.75em]

  \caption{Hierarchical variance allocation for the Beta-Beta model,
    $\dpm = 0.5$, $\dpv = 0.04$.}
  \label{fig:bb_sensitivity}
\end{figure}

\begin{figure}[htbp]
  \centering
  \small (a) KS distance: $\sup_z |F_{\text{compound}}(z) - F_{\text{direct}}(z)|$. \\
  \includegraphics[width=0.6\textwidth]{../plots/paper/bb_ks_distance.pdf}\\[-.75em]

  \vspace{0.8em}
  \small (b) CDF difference at four $(\rho, \kappa)$ points. \\
  \includegraphics[width=0.8\textwidth]{../plots/paper/bb_cdf_diff_curves.pdf}\\[-.75em]

  \caption{Distributional approximation quality for the Beta-Beta hierarchy,
    $\dpm = 0.5$, $\dpv = 0.04$.}
  \label{fig:bb_distributional}
\end{figure}

Figures~\ref{fig:bb_sensitivity} and \ref{fig:bb_distributional} provide the
graphical calibration for the Beta-Beta hierarchy with $\dpm = 0.5$ and $\dpv =
0.04$. The sensitivity heatmaps show qualitatively similar patterns to the NIG
case, with the notable difference that regions of $(\rho, \kappa)$ space where
the Beta-Beta hyperprior is infeasible appear as missing values. The KS
distances (panel~(a) of Figure~\ref{fig:bb_distributional}) are substantially
larger than the NIG case, reaching 0.20 at the extreme feasible corner compared
to 0.04 for NIG. This reflects the additional distributional distortion
introduced by the cross-moment coupling between $\mu_k$ and $\psi_k$: even
after moment-matching, the Beta-Beta compound departs further from the direct
prior than the NIG compound does.

%----------------------------------------------------------------------
\subsection{Summary}
\label{app:family_summary}

The three families illustrate the range of behavior that arises when translating
$(\rho, \kappa)$ to hyperprior parameters. The NIG and Gamma-IG families share
the simplifying property that the conditional variance is a single parameter
($\sigma^2$) with its own independent hyperprior, making $\kappa$ map directly
to the Inverse-Gamma shape parameter. Both admit the full $(\rho, \kappa)$
region with no joint constraints. The Beta-Beta family, by contrast, has a
conditional variance that depends on both hyperparameters, requiring a
cross-moment correction and imposing a lower bound on $\kappa$ that increases
with $\rho$.

The distributional approximation quality follows a similar ordering. The NIG
compound achieves the smallest KS distances (below 0.05 across the entire
grid), the Gamma-IG is intermediate (0.005 to 0.14), and the Beta-Beta is
largest (up to 0.20 at extreme feasible values). In all three cases, the
moment-matching construction guarantees that the first two moments of the
compound distribution match the direct prior exactly; only the higher moments
differ, with the discrepancy concentrated in the tails and shoulders of the
distribution.
